% c06 —— §7 末（Hesse 群与 X_0(3)）+ §8 与 Q(sqrt5) 的对称 Hilbert 模曲面相伴的 K3 曲面
%        + §8.1 二次栈（p014–p016 / 原书 p.140–142）
% 承 c05：原书 p.139 末句「Hesse 群 $G_{216}\cong\cdots$ は $S(3)\to X(3)$ に同変に作用する」跨页，
% 整句归本块，从「Hesse 群 …… 同变地作用于 $S(3)\to X(3)$。」起写。
% 原书 p.141 的首句「群で割ったものになる．」承接 p.140 末的「模曲线直积……」，故 p.141 的
% 页码标记放在其后的「二次栈给出族参数空间的紧化……」段之前。
% 结尾止于原书 p.142 的 (8.10)（判别式长式）；p.143 开头「であり，これは斉次座標では」归 c07。

% ===== 原书 p.140 =====

Hesse 群（Hessian group）$G_{216}\cong SL_{2}(\mathbb{F}_{3})\ltimes(\mathbb{Z}/3\mathbb{Z})^{2}$ 同变地作用于
$S(3)\to X(3)$。这里 $SL_{2}(\mathbb{F}_{3})=SL_{2}(\mathbb{Z})/\Gamma(3)$，且
$PSL_{2}(\mathbb{F}_{3})\cong A_{4}$。设 $\omega=\exp(2\pi\sqrt{-1}/3)$ 为 $1$ 的本原 $3$ 次单位根，
考虑由
\begin{equation}
  ([X:Y:Z],\mu)\mapsto([\omega X:\omega^{2}Y:Z],\mu), \tag{7.6}
\end{equation}
\begin{equation}
  ([X:Y:Z],\mu)\mapsto([X:Y:\omega Z],\omega\mu), \tag{7.7}
\end{equation}
生成的 $G_{216}$ 的子群。对每个 $\mu$，由 (7.3) 的特殊化得到的椭圆曲线
\begin{equation}
  x^{2}y+xy^{2}+z^{3}+\mu xyz=0 \tag{7.8}
\end{equation}
是 (7.5) 按 (7.6) 取商所得的：
\begin{equation}
  x=\frac{X^{2}}{Y},\qquad y=\frac{Y^{2}}{X},\qquad z=Z,\qquad \nu=\frac{1}{\mu^{3}}. \tag{7.9}
\end{equation}
% [Tutor Note: (7.9) 末项印作 $\nu=1/\mu^{3}$，而 (7.8) 中的参数印作 $\mu$，两者在记号上不衔接；此处按原图照录，未作统一。]
在此映射下，每个纤维上的子群 $E[3]$ 被映到 $E/(\mathbb{Z}/3\mathbb{Z})$ 的与
$\mathbb{Z}/3\mathbb{Z}$ 同构的子群上。底空间上的 2 个点 $\mu$ 与 $\omega\mu$
作为带 $\Gamma_{0}(3)$ 结构的椭圆曲线彼此同构，取该作用的商，便得到椭圆模曲面
$S_{0}(3)\to X_{0}(3)$。模曲线 $X_{0}(3)$ 有 2 个尖点，以及 1 个具有 3 阶稳定化子的点。

\section{与 $\mathbb{Q}(\sqrt{5})$ 的对称 Hilbert 模曲面相伴的 K3 曲面}

3 维自反多面体共有 4319 个，其中顶点个数为 5 的（即对应的 K3 曲面族的参数个数为 2 的）有 5 个。
[HNU] 从二次栈与环面紧化的观点详细研究了其中 2 个。
这些族的一方是 $\mathbb{P}^{2}$ 上 $\mathbb{P}^{1}$ 丛
$\mathbb{P}(\mathcal{O}_{\mathbb{P}^{2}}\oplus\mathcal{O}_{\mathbb{P}^{2}}(2))$ 的反典范超曲面的镜像，
另一方是 Picard 数为 1、具有 1 个通常二重点的环面 Fano 流形的小消解的反典范超曲面的镜像。
这些族的非常一般的纤维的 Picard 格与超越格分别为
\begin{equation}
  M=E_{8}\perp E_{8}\perp\begin{pmatrix}2&1\\1&-2\end{pmatrix},\qquad
  T=U\perp\begin{pmatrix}2&1\\1&-2\end{pmatrix} \tag{8.1}
\end{equation}
以及
\begin{equation}
  M=E_{8}\perp E_{8}\perp\begin{pmatrix}0&3\\3&0\end{pmatrix},\qquad
  T=U\perp\begin{pmatrix}0&3\\3&0\end{pmatrix} \tag{8.2}
\end{equation}
[Nag12]。对应的格极化 K3 曲面的模空间 $\mathcal{M}$，前者成为与 $\mathbb{Q}(\sqrt{5})$
相伴的对称 Hilbert 模曲面，后者则成为模曲线的直积用适当群取商所得的商。

% ===== 原书 p.141 =====

二次栈给出族参数空间的紧化，它与模空间的环面紧化的关系是 [HNU] 的主定理；
在 Hilbert 模曲面的情形，该关系如下。

\begin{theorem}[{[HNU, Theorem 1.1]}]
周期映射给出 Deligne–Mumford 栈的同构
\begin{equation}
  \tilde{\Pi}\colon\tilde{X}_{\mathcal{F}(A)}\xrightarrow{\sim}\mathcal{M}_{\Sigma}. \tag{8.3}
\end{equation}
这里 $\tilde{X}_{\mathcal{F}(A)}\to X_{\mathcal{F}(A)}$ 是在 $X_{\mathcal{F}(A)}$ 的一点处按权 $(1,3)$
作加权爆破之后，再沿判别式作二阶根构造所得的。又，$\mathcal{M}_{\Sigma}$ 是在由周期映射的
单值性确定的扇 $\Sigma$ 所给出的 $\mathcal{M}$ 的环面紧化上赋予自然的轨形结构所得的。
\end{theorem}

这里，根构造（root construction）是沿除子引入稳定化子的方法 [AGV08, Cad07]。
佐竹–Baily–Borel 紧化 $\overline{\mathcal{M}}_{0}$ 的唯一尖点的原像是 2 个环面除子的并，
其交点是极大退化点（maximally unipotent monodromy point）。
这些除子附近的单值性对数 $N$ 满足 $N^{2}\neq0$ 与 $N^{3}=0$，由此可知 K3 曲面发生 III 型退化。
把周期映射表述为有理映射，这一点已在 [Nag12] 中得到，而定理 8.1 进一步
\begin{itemize}[leftmargin=2.2em]
  \item 用二次栈 [Laf03, Hac, DKK] 紧化参数空间，
  \item 用 [Iri11] 把周期映射的单值性用镜像的相干层的导出范畴的自等价的语言描述，
  \item 用该单值性作模空间的环面紧化，
  \item 把周期映射描述为紧栈之间的映射。
\end{itemize}

定理 8.1 的动机在于基于对数几何的 K3 曲面模空间紧化 [Ols04] 与对数 Hodge 理论 [KU09]，
以及对数几何在镜像对称中日益提高的重要性（例如参见 [ACG$^{+}$, Gro]，以及本《講究録》中臼井氏的论说）。

\subsection{二次栈}

设图 8.1 中 5 顶点的自反多面体为 $\Delta=\Conv\{v_{1},v_{2},v_{3},v_{4},v_{5}\}$，则
\begin{equation}
  \beta=(v_{1}\ v_{2}\ v_{3}\ v_{4}\ v_{5})=
  \begin{pmatrix}
    1&0&0&0&-1\\
    0&1&0&0&-1\\
    0&0&1&-1&-2
  \end{pmatrix}. \tag{8.4}
\end{equation}
% [Tutor Note: 原书图 8.1 中各点标作 $a_{0},\dots,a_{5}$（$a_{0}$ 疑为内点），正文与 (8.4) 却用 $v_{1},\dots,v_{5}$；此处照原文分别保留。]

% ===== 原书 p.142 =====

\begin{figure}[htbp]
  \centering
  \includegraphics[width=0.40\textwidth]{fig8_1.png}
  \caption{5 顶点的自反多面体 $\Delta$（原书 p.142）}
\end{figure}

把 (5.4) 的除子序列中出现的同态 $\mathbb{Z}^{n+r}\to\LL^{\vee}$ 写成矩阵，则得
\begin{equation}
  \begin{pmatrix}
    -2&0&0&1&1&0\\
    -5&1&1&2&0&1
  \end{pmatrix}. \tag{8.5}
\end{equation}
由此二次扇如图 8.2 所示。

\begin{figure}[htbp]
  \centering
  \includegraphics[width=0.40\textwidth]{fig8_2.png}
  \caption{二次扇 $\mathcal{F}(A)$（原书 p.142）}
\end{figure}

Laurent 多项式空间由
\begin{equation}
  \mathbb{C}^{A_{0}}=\Bigl\{W=a_{0}+a_{1}x+a_{2}y+a_{3}z+\frac{a_{4}}{z}+\frac{a_{5}}{xyz^{2}}\Bigr\} \tag{8.6}
\end{equation}
给出。二次栈 $X_{\mathcal{F}(A_{0})}$ 的稠密环面 $\LL^{\vee}_{\mathbb{C}^{\times}}$ 可用
\begin{equation}
  \lambda=\frac{a_{3}a_{4}}{a_{0}^{2}}, \tag{8.7}
\end{equation}
\begin{equation}
  \mu=\frac{a_{1}a_{2}a_{3}^{2}a_{5}}{a_{0}^{5}} \tag{8.8}
\end{equation}
写成 $\Spec\mathbb{C}[\lambda^{\pm1},\mu^{\pm1}]$。二次栈 $X_{\mathcal{F}(A)}$ 由对
$\mathbb{P}(1,2,5)$ 作加权爆破得到，爆破前 $\mathbb{P}(1,2,5)$ 的齐次坐标由
\begin{equation}
  [\nu:\lambda:\mu]=[a_{0}:a_{3}a_{4}:a_{1}a_{2}a_{3}^{2}a_{5}] \tag{8.9}
\end{equation}
给出。判别式为
\begin{align}
  \Delta={}&a_{4}^{2}a_{0}^{6}+4a_{1}a_{2}a_{5}a_{0}^{5}-12a_{3}a_{4}^{3}a_{0}^{4}
    -50a_{1}a_{2}a_{3}a_{4}a_{5}a_{0}^{3}\notag\\
  &+48a_{3}^{2}a_{4}^{4}a_{0}^{2}+1000a_{1}a_{2}a_{3}^{2}a_{4}^{2}a_{5}a_{0}
    -64a_{3}^{3}a_{4}^{5}+3125a_{1}^{2}a_{2}^{2}a_{3}^{2}a_{5}^{2}. \tag{8.10}
\end{align}
% [Tutor Note: 原书 (8.10) 末行以句点结束，而原书 p.143 开头的「であり，これは斉次座標では」在语法上承接 (8.10)，该句点疑为逗号误排；此处按原图保留句点。]

% CHUNK c06 DONE pages p014–p016
